% Transcript-derived local integration, 2026-09-16.
% Source transcript SHA256 e118ecee57ea4d6b93489f62ddaedcc83e63165399b12c064b5daae4ba9338bc.
% Accepted file14 SHA256 45cf9435afdea364495331da06272622ed356257c155f377c19def812af6450a.
% Additive body: requires amsmath, amssymb. No predecessor is edited.
\section{The complete original action relative to its arithmetic total}
\label{sec:complete-action-integration}

This is a local reconstruction of the complete delivered response following
the 16 September 2026, 15:02:50 prompt in
\texttt{ChatGPT-Residual Lower Bound Continuation-20260916-2150.md},
lines 7043--7730. The response reported that its build tools timed out:
no original standalone archive, compiled PDF, successful fresh Lean run or
actual-period numerical certificate is inferred from it. This section
supplies the additive proof in full. Its input source is the unchanged
\texttt{14\_CURRENT\_KERNEL\_AND\_MULTIPLICITY\_PROOFS.tex}; the original
AMT, MCE and OMR formulas below are rederived at the points where used.
The fixed arithmetic density-envelope constants are precisely AMT4--6's
proved constants. Their defining bounds and domains are displayed, rather
than replaced by newly selected measures or a new positivity premise.

\subsection{Domain, unchanged arithmetic objects and exact minimum metrics}

Use the original stipulated full zero-quartet domain: fixed
$0<\delta<1/2$, $\gamma>2$, and complete common order $m\ge1$ at the four
distinct centres. The statement below concerns these original objects; it
does not assert that such an off-line zero quartet exists. Write
\[
\begin{gathered}
 \rho=\tfrac12+\delta+i\gamma,\quad
 \mathcal R=\{\rho,\bar\rho,1-\rho,1-\bar\rho\},\quad
 h(s)=\prod_{\zeta\in\mathcal R}(s-\zeta)^m,\quad g=2\xi,\quad v_h=g/h,\\
 w_h(y)=|v_h(1/2+iy)|^2/(2\pi),\quad m_{h,b}=w_h^{*b},\quad
 \mu_h=\int_{\mathbb R}w_h(y)\,dy,\\
 k=4\ell+1\ge9,\quad c=k/2,\quad e_k=1+k(m-1),\quad
 q=e_k(k+1)^2,\\
 \beta_{ab}=c+(2a-k)\delta+i(2b-k)\gamma,\quad
 \chi(S)=\prod_{a,b=0}^{k}(S-\beta_{ab})^{e_k},\quad
 E=\mathbb C[S]/(\chi),\quad A=M_S.
\end{gathered}\tag{CAI1}
\]
The arithmetic source is $L^2(\mathbb R,m_{h,k}(y)\,dy)$ with its
mass $\mu_h^k$, evaluated in the literal coordinate $S=c+iy$.
Its inner product is conjugate-linear in the first argument. For each
$N\ge0$ let $\mathcal P_N=\mathbb C[S]_{\le N}$ and let
$p_n(S)=i^nQ_n((S-c)/i)$ be the monic orthogonal polynomials; $Q_n$ is
monic in the real variable $y$. Their full squared norms are
$\omega_{n,k}^{\rm ar}>0$. The substitution has inverse
$P(S)=f((S-c)/i)$ and is never a change of the original source mass.
In formulas below $\omega_n$ abbreviates this same arithmetic norm.

The exact remainder map $J_N:\mathcal P_N\to E$ is onto for
$N\ge q-1$, with kernel $\chi\mathcal P_{N-q}$, where
$\mathcal P_{-1}=0$. Put $b_n=[p_n]_\chi$ in the frame
$1,S,\ldots,S^{q-1}$. Then
\[
 R_N=\sum_{n=0}^N b_nb_n^*/\omega_n,\qquad
 G_N=R_N^{-1},\qquad V_N=\det G_N.\tag{CAI2}
\]
The first $q$ columns form a unit-triangular monic remainder matrix, so
$R_N>0$. In any source coefficient frame with Gram $H_N$, the actual
minimum section is $H_N^{-1}J_N^*G_N$. Its composition with $J_N$ is the
identity. A vector in $\ker J_N$ pairs to zero with this section;
therefore the squared norm of another lift is its minimum squared norm
plus the squared norm of that relation. This proves the interpretation
of $G_N$ without changing the admitted relation subspace.

The original period observation is the AMT3 map
$\Lambda:E\twoheadrightarrow B$, with actual kernel inclusion
$I:K\hookrightarrow E$. Explicitly, if $E_h=\mathbb C[s]/(h)$,
$U=M_{j_hv_h}$, $\iota[P]=[P(s_1+\cdots+s_k)]$, $\Pi$ is the original
period matrix and $P_k$ its original tensor invariant projector, then
$\jmath\Lambda=P_k\Pi^{\otimes k}U^{\otimes k}\iota$, where
$\jmath:B\hookrightarrow(\mathbb C^{4m})^{\otimes k}$ is the image
inclusion. No period branch or kernel frame is changed. The action
calculation below is independent of that choice because it uses the
complete quotient $E$; the exact relation to the original observation
is the following splitting. In the fixed kernel and image frames set
\[
 H_{K,N}=I^*G_NI,\quad
 Q_{B,N}=(\Lambda G_N^{-1}\Lambda^*)^{-1},\quad
 S_N=G_N^{-1}\Lambda^*Q_{B,N}.\tag{CAI3}
\]
Direct multiplication proves $\Lambda S_N=1$ and $I^*G_NS_N=0$.
For any fixed section $S_*$ of $\Lambda$, the frame $[I,S_*]$ is
invertible and $S_N-S_*$ has image in $K$. Thus changing $[I,S_*]$ to
$[I,S_N]$ has determinant one. Orthogonality proves the exact identity
\[
 \det H_{K,N}\det Q_{B,N}=|\det[I,S_*]|^2V_N.\tag{CAI4}
\]
Empty factors have determinant one, so this also covers zero-dimensional
kernel or image. In particular no kernel invariance under $A$ is needed.

For any one of the original sources $a={\rm ar},1,k$, define its full
two-window return and the kernel and boundary returns by
\[
\begin{gathered}
 \mathcal B_k^a=\log\frac{V_{q-1}^aV_q^a}{V_{2q-1}^aV_{2q}^a},\\
 F_K^a=\log\frac{\det H_{K,q-1}^a\det H_{K,q}^a}
                       {\det H_{K,2q-1}^a\det H_{K,2q}^a},\quad
 F_B^a=\log\frac{\det Q_{B,q-1}^a\det Q_{B,q}^a}
                       {\det Q_{B,2q-1}^a\det Q_{B,2q}^a},\\
 F_K^a+F_B^a=\mathcal B_k^a,\qquad
 \delta_k^\sigma=\mathcal B_k^{\rm ar}-\mathcal B_k^\sigma,\quad
 \mathcal H_k=\mathcal B_k^\sigma-\mathcal B_k^{(k)}.
\end{gathered}\tag{CAI5}
\]
Here $a=1$ denotes the original fixed Gamma source $\sigma$.
The equality follows by multiplying CAI4 at the two lower degrees and
dividing by its two upper-degree instances. Its four identical frame
constants cancel exactly, rather than being set equal to one.

\subsection{The actual rank-two control, reflection and exterior benchmark}

The quotient $v_h$ is entire at the selected centres because their
complete zero orders have been divided out. The exact equalities
$\xi(1-s)=\xi(s)$ and $h(1-s)=h(s)$ show that $w_h$ and every $m_{h,b}$
are even. Monic orthogonality consequently gives $Q_n(-y)=(-1)^nQ_n(y)$.
Self-adjoint multiplication by $y$, orthogonality to degrees below
$n-1$, and pairing with $Q_{n-1}$ prove the real monic recurrence.
In the original $S$ coordinate it is
\[
 Sp_n=p_{n+1}+cp_n-(\omega_n/\omega_{n-1})p_{n-1},\tag{CAI6}
\]
with the lower term absent at $n=0$. Inserting its quotient columns into
CAI2 cancels every interior adjacent pair and gives
\[
\begin{gathered}
 AR_N+R_NA^*-kR_N=(b_{N+1}b_N^*+b_Nb_{N+1}^*)/\omega_N,\\
 D_N^{\rm ctrl}=A^*G_N+G_NA-kG_N,\quad
 G_N^{-1/2}D_N^{\rm ctrl}G_N^{-1/2}=(yx^*+xy^*)/\omega_N,\\
 x=G_N^{1/2}b_N,\qquad y=G_N^{1/2}b_{N+1}.
\end{gathered}\tag{CAI7}
\]
The symbol $D_N^{\rm ctrl}$ is distinct from the scalar source constant
$D_N$ below. Reflection $P(S)\mapsto P(k-S)$ preserves $\chi$ because
it permutes $(a,b)$ with $(k-a,k-b)$ and $q$ is even. It therefore induces
an involution $\Sigma$ on $E$, with $\Sigma A\Sigma=kI-A$.
Changing variables $y\mapsto-y$ is an isometry of each entire remainder
fibre. Uniqueness of the minimum proves $\Sigma^*G_N\Sigma=G_N$.
Parity gives $\Sigma b_n=(-1)^nb_n$, and hence
$b_N^*G_Nb_{N+1}=-b_N^*G_Nb_{N+1}=0$. Thus the only possibly nonzero
eigenvalues of the relative Hermitian control in CAI7 are
$\epsilon_N,-\epsilon_N$, where
\[
 \epsilon_N^2=\frac{\|x\|^2\|y\|^2}{\omega_N^2},\qquad
 \phi_N=0.\tag{CAI8}
\]
For comparison with the original tilted formulas, before specialization
the exact expression is
$(\|x\|^2\|y\|^2-|x^*y|^2)/\omega_N^2$ and
$x^*y=i\omega_N\phi_N$. At the present untilted source only the
\emph{total} pairing vanishes. In the orthogonal kernel--boundary split
write $x=x_K\oplus x_B$, $y=y_K\oplus y_B$, and
$z_K=x_K^*y_K$, $z_B=x_B^*y_B$. Then $z_B=-z_K$, not generally zero.
In particular the mixed exterior summand remains
\[
 \|x_K\otimes y_B-y_K\otimes x_B\|^2
 =\|x_K\|^2\|y_B\|^2+\|x_B\|^2\|y_K\|^2+2|z_K|^2.\tag{CAI9}
\]

Here is the original source-to-class map and its lower allowance. The
complete Taylor unit $\upsilon_h=j_hv_h$ is invertible in $E_h$, and
\[
 \eta:E\longrightarrow E_h^{\otimes k},\qquad
 \eta[P]=\upsilon_h^{\otimes k}P(A_\Sigma)1,\qquad
 A_\Sigma=\sum_iM_{s_i},\qquad \eta A=A_\Sigma\eta.\tag{CAI10}
\]
In each local tensor factor the nilpotent has order $m$. The sum of the
$k$ nilpotents has exact order $e_k$: the power $e_k-1=k(m-1)$ has
nonzero top coefficient $[k(m-1)]!/[(m-1)!]^k$, and every term of its
next power has one exponent at least $m$. Every centre $\beta_{ab}$
occurs, and distinct pairs have distinct centres. The annihilator of
$A_\Sigma$ on its cyclic unit is therefore exactly $\chi$; this proves
injectivity of CAI10, including all jets. Reflection fixes
$\upsilon_h$ and intertwines this inclusion with the original tensor
reflection $s_i\mapsto1-s_i$.

Let $E_+$ be the direct sum of all complete primary blocks with $a>k/2$.
Its dimension is $t_+=q/2$. A primary idempotent is explicitly
$\chi_\beta\sum_{j=0}^{e_k-1}(1/\chi_\beta)^{(j)}(\beta)
(S-\beta)^j/j!\pmod\chi$, where $\chi_\beta=\chi/(S-\beta)^{e_k}$.
Its product with the powers $(S-\beta)^d$, $0\le d<e_k$, gives a
complete primary basis. The top wedge $\xi_+$ of these bases is nonzero,
and remains nonzero under $\bigwedge^{t_+}\eta$. Its additive action
eigenvalue is $\operatorname{Tr}(A|_{E_+})$; nilpotents contribute zero
to that trace. The relative Hermitian additive action on any exterior
degree $1\le t\le q-1$ has largest eigenvalue $\epsilon_N$, since its
eigenvalues are sums of distinct members of
$\{\epsilon_N,-\epsilon_N,0,\ldots,0\}$. Its expectation on $\xi_+$
is $2\operatorname{Re}\operatorname{Tr}(A|_{E_+})-kt_+$. Summing the
positive odd numbers $2a-k$ gives $(k+1)^2/4$, so
\[
 0<L_{h,k}:=2\delta e_k(k+1)\frac{(k+1)^2}{4}
          =\frac{\delta q(k+1)}2\le\epsilon_N
          \qquad(N\ge q-1).\tag{CAI11}
\]
This is a calculation on the same stipulated arithmetic class, not a
claim of its existence and not an assumed desired upper estimate.

\subsection{Exact cancellation before any estimates}

Set $d_n=V_n/V_{n-1}$ for $n\ge q$, and
$b_n^{\rm rec}=\omega_n/\omega_{n-1}$; neither is the centre displacement
$\delta$ or the quotient column $b_n$. From the rank-one covariance
update and the determinant lemma,
\[
 \frac{b_N^*G_Nb_N}{\omega_N}=1-d_N\quad(N\ge q),\qquad
 \frac{b_{N+1}^*G_Nb_{N+1}}{\omega_{N+1}}=d_{N+1}^{-1}-1.
\]
The first equality follows by removing the last covariance column;
the second follows by adding the next one. At $N=q-1$ the first $q$
columns are the full quotient basis and
$b_{q-1}^*G_{q-1}b_{q-1}=\omega_{q-1}$, without using any inverse at
$q-2$. Thus CAI8 gives the original MCE23/OMR18 formulas
\[
 \epsilon_{q-1}^2=b_q^{\rm rec}(d_q^{-1}-1),\qquad
 \epsilon_N^2=b_{N+1}^{\rm rec}(1-d_N)(d_{N+1}^{-1}-1) (N\ge q).
 \tag{CAI12}
\]
Positive covariances give $0<d_n\le1$, and CAI11--12 force $d_n<1$
at every degree appearing here. All following logarithms therefore
exist. Put $w_{q-1}=w_{2q-1}=1$ and $w_N=2$ for $q\le N\le2q-2$;
their sum is $2q$. Retain exactly
\[
\begin{gathered}
 \mathcal W_k^{\rm ar}=\log\frac{\omega_{2q-1}\omega_{2q}}
                                      {\omega_{q-1}\omega_q},\\
 \mathcal P_-=-\log(1-d_{2q-1})-2\sum_{n=q}^{2q-2}\log(1-d_n),\\
 \mathcal P_+=-\log(1-d_q)-\log(1-d_{2q})
                              -2\sum_{n=q+1}^{2q-1}\log(1-d_n),\\
 J_k^{\rm action}=2\sum_{N=q-1}^{2q-1}w_N\log\epsilon_N,\quad
 \Psi_k=F_K^{\rm ar}+F_B^{\rm ar}-F_B^{(k)},\quad
 C_k^{\rm budget}=F_B^{(k)}+\mathcal W_k^{\rm ar}-\mathcal P_--\mathcal P_+.
\end{gathered}\tag{CAI13}
\]
Multiply CAI12 over $N=q-1,\ldots,2q-2$ and over
$N=q,\ldots,2q-1$ separately. The two exact products are
\[
\begin{aligned}
 \prod_{N=q-1}^{2q-2}\epsilon_N^2
 &=\frac{\omega_{2q-1}V_{q-1}}{\omega_{q-1}V_{2q-1}}
            (1-d_{2q-1})\prod_{n=q}^{2q-2}(1-d_n)^2,\\
 \prod_{N=q}^{2q-1}\epsilon_N^2
 &=\frac{\omega_{2q}V_q}{\omega_qV_{2q}}
            (1-d_q)(1-d_{2q})\prod_{n=q+1}^{2q-1}(1-d_n)^2.
\end{aligned}\tag{CAI14}
\]
Indeed $d_{N+1}^{-1}-1=(1-d_{N+1})/d_{N+1}$; norm ratios telescope,
reciprocal $d$ factors give the displayed volume ratios, and the
endpoint and interior contraction exponents are precisely those written.
Taking logarithms and using CAI5 proves
\[
 \boxed{J_k^{\rm action}=\Psi_k+C_k^{\rm budget}
    =\mathcal B_k^{\rm ar}+\mathcal W_k^{\rm ar}-\mathcal P_--\mathcal P_+.}
 \tag{CAI15}
\]
Equivalently $\Psi_k=F_K^{(k)}+\mathcal H_k+\delta_k^\sigma$ and
$C_k^{\rm budget}=\mathcal B_k^{(k)}-F_K^{(k)}+
\mathcal W_k^{\rm ar}-\mathcal P_--\mathcal P_+$.
The cancellation uses the exact joint identities, not independent
estimates for the kernel, residual, low metrics or transfer errors.

\subsection{All-degree multiplication bounds on the actual arithmetic blocks}

The AMT1 domain is every convolution order $b\ge3$, not just the later
tensor subsequence. Its constants, without alteration, are
\[
\begin{gathered}
 \alpha=\pi/2,\quad p=8m-\tfrac92,\quad B_h=42+8m,\quad M=21+4m,\\
 \vartheta_h=\int_{-1}^1w_h(y)\,dy>0,\quad
 M_\alpha=\int_{\mathbb R}e^{\alpha y}w_h(y)\,dy\in(0,\infty),\\
 c_\Gamma=\sqrt2e^{-7/6},\quad C_\Gamma=\sqrt{2\sqrt5}e^{2/3},\\
 a_b=\frac{c_h\vartheta_h^{b-3}e^{-\alpha(b-3)}(b-2)^{-B_h}}
                   {C_\Gamma2^{B_h/2}},\quad
 A_b=\frac{C_h^{\rm tilt}}{c_\Gamma}b^{p+1}M_\alpha^{b-1},\quad
 D_N=[1+4(N+M)^2]^M.
\end{gathered}\tag{CAI16}
\]
Here $c_h$ is the original TW7 lower-envelope constant and
$C_h^{\rm tilt}$ the original BT12 compact/tail maximum, as recorded
in AMT4. Their already proved defining inequalities, for every $b\ge3$,
are
\[
\begin{aligned}
 m_{h,b}(y)&\ge c_h\vartheta_h^{b-3}
             e^{-\alpha(|y|+b-3)}(b-2+|y|)^{-B_h},\\
 m_{h,b}(y)&\le bC_h^{\rm tilt}M_\alpha^{b-1}
             e^{-\alpha|y|}(1+|y|/b)^{-p},\\
 c_\Gamma e^{-\alpha|y|}(1+|y|)^{-1/2}
 &\le \frac{d\sigma}{dy}\le
 C_\Gamma e^{-\alpha|y|}(1+|y|)^{-1/2},\qquad
 d\sigma=|\Gamma(1/4+iy/2)|^2dy/(2\pi).
\end{aligned}\tag{CAI17}
\]
The factors $b-2+|y|\le(b-2)(1+|y|)$,
$(1+|y|)^{B_h}\le2^{B_h/2}(1+y^2)^M$, and
$1+|y|/b\ge(1+|y|)/b$, with $p>1/2$, give
$a_b(1+y^2)^{-M}d\sigma/dy\le m_{h,b}\le A_b d\sigma/dy$.
The orthonormal Gamma recurrence is
$ye_j=\sqrt{(j+1)(j+1/2)}e_{j+1}+\sqrt{j(j-1/2)}e_{j-1}$.
It gives $\|yP\|_\sigma\le2(N+1)\|P\|_\sigma$ on degree $N$.
Iterate at the successive degrees and expand $(1+y^2)^M$ to get
$\int|P|^2(1+y^2)^M d\sigma\le D_N\|P\|_\sigma^2$.
Cauchy--Schwarz applied with the two weights $(1+y^2)^{\pm M/2}$
therefore proves, for every complex polynomial of degree at most $N$,
\[
 \frac{a_b}{D_N}\|P\|_\sigma^2\le\|P\|_{m_{h,b}}^2
                         \le A_b\|P\|_\sigma^2.\tag{CAI18}
\]
These are precisely the full-polynomial AMT6 inequalities; in particular
they apply at $b=3,4,5$ with no new source-order restriction.

The fixed Gamma monic polynomials have full norms
$\gamma_n=\sqrt{2\pi}\,n!(1/2)_n$. Their recurrence gives the
generating series $F(y,z)=\sum_{n\ge0}p_n^\sigma(y)z^n/n!$:
\[
 (1+z^2)\partial_zF=(y-z/2)F,\quad F(y,0)=1,\qquad
 F=(1+z^2)^{-1/4}e^{y\arctan z}.\tag{CAI19}
\]
One may verify the differential equation coefficient by coefficient,
and integration at zero uniquely gives the displayed analytic function
on $|z|<1$. On $|z|=1/4$, its prefactor has modulus at most
$(16/15)^{1/4}$ and
$|\arctan z|\le\sum_{j\ge0}|z|^{2j+1}/(2j+1)=\operatorname{arctanh}(1/4)$.
Consequently Cauchy's coefficient bound and
$n!/(1/2)_n=4^n/\binom{2n}{n}\le2n+1$ imply
\[
 |P(y)|^2\le C_\sigma(N+1)(2N+1)16^N e^{\nu_0|y|}\|P\|_\sigma^2,
 \quad C_\sigma=\frac4{\sqrt{15}\sqrt{2\pi}},\quad
 \nu_0=\log(5/3).\tag{CAI20}
\]
For the binomial inequality the central coefficient is the maximum of
the $2n+1$ positive coefficients summing to $4^n$. For the polynomial
inequality expand in the complete degree-$N$ orthonormal basis and
apply Cauchy--Schwarz to the coefficient vector; then bound each of
the $N+1$ summands by $(2N+1)16^N$. This covers complex coefficients.

Put $d_0=\alpha-\nu_0>0$. Positivity and convolution give
$\int e^{\alpha y}m_{h,b}(y)\,dy=M_\alpha^b$ by Tonelli, and evenness
gives $\int e^{\alpha|y|}m_{h,b}(y)\,dy\le2M_\alpha^b$.
For $x>R\ge0$,
$x^2e^{-d_0x}\le(16/d_0^2)e^{-d_0R/2}$: the maximum of
$x^2e^{-d_0x/2}$ is $16/(e^2d_0^2)\le16/d_0^2$.
Moreover $N+M\le M(N+1)$ since $M\ge1$, so
$D_N\le(1+4M^2)^M(N+1)^{2M}$. Combining these facts with CAI18--20,
and $2N+1\le2(N+1)$, proves for $b=3,4,5$
\[
 \int_{|y|>R}y^2|P(y)|^2m_{h,b}(y)\,dy
 \le H_h(N+1)^{2M+2}16^Ne^{-d_0R/2}\|P\|_{m_{h,b}}^2,
 \quad H_h=\frac{64C_\sigma(1+4M^2)^M}{d_0^2}
                          \max_{b=3,4,5}\frac{M_\alpha^b}{a_b}.
 \tag{CAI21}
\]
All constants are fixed-$h$ data, including the unscaled source masses.
Choose
\[
 T_h=\frac2{d_0}\{\log16+2M+3+\log^+H_h\},\qquad
 C_{\rm mult}(h)=\sqrt{T_h^2+1}.\tag{CAI22}
\]
Here $\log^+x=\max(0,\log x)$. At $R=T_h(N+1)$ the logarithm of the
factor in CAI21 is at most
$\log H_h+(2M+2)N+N\log16-
(\log16+2M+3+\log^+H_h)(N+1)\le0$, using
$\log(N+1)\le N$. The inner interval has $y^2\le T_h^2(N+1)^2$.
Adding its integral and the tail, and using $(N+1)^2\ge1$, proves
\[
 \|yP\|_{m_{h,b}}\le C_{\rm mult}(h)(N+1)\|P\|_{m_{h,b}},
                         \qquad b=3,4,5.\tag{CAI23}
\]

\subsection{Total-degree tensor compression, including the output face}

For every integer $k\ge3$, write $r_k=\lfloor k/3\rfloor$ and
partition $k$ into $r_k$ blocks of sizes $3,4,5$. If $k=3r_k$ take
all threes; if $k=3r_k+1$ or $3r_k+2$, replace one three by four or
five. This also covers $k=3,4,5$. Convolution proves the exact isometry
\[
 L^2(m_{h,k})\longrightarrow\bigotimes_{i=1}^{r_k}L^2(m_{h,b_i}),
 \quad f(y)\longmapsto f(y_1+\cdots+y_{r_k}),\quad
 \sum_i(b_i/2+iy_i)=k/2+i\sum_iy_i.\tag{CAI24}
\]
The norm identity follows first by Tonelli for $|f|^2$ and then by
the definition of convolution. Its image is the sum-dependent
subspace, not the whole product space; total masses multiply to
$\mu_h^k$. In each block let $\varphi_n^{(b)}$ be the real orthonormal
polynomial with positive leading coefficient. Its exact recurrence is
\[
 y\varphi_n^{(b)}=a_{n+1}^{(b)}\varphi_{n+1}^{(b)}
                         +a_n^{(b)}\varphi_{n-1}^{(b)},\qquad
 a_0^{(b)}=0,\quad 0<a_n^{(b)}\le C_{\rm mult}(h)n\ (n\ge1).
\]
There is no diagonal term by evenness. The bound follows by taking
the coefficient of $\varphi_n$ in $y\varphi_{n-1}$ and applying CAI23.

Use all product basis vectors of total degree at most $N+1$, and let
$M_{N+1}$ be the compression of multiplication by $\sum_i y_i$ to
their span. It is a real symmetric matrix. At a vertex
$\nu=(\nu_1,\ldots,\nu_{r_k})$ of total degree at most $N$, its full
absolute row sum is at most
$C_{\rm mult}\sum_i(\nu_i+(\nu_i+1))\le C_{\rm mult}(2N+r_k)$.
At total degree $N+1$ only downward entries remain in the compression,
so that row sum is at most $C_{\rm mult}(N+1)\le C_{\rm mult}(2N+r_k)$.
Symmetry gives the same bound for column sums and hence the operator
norm: directly, for any vector $z$, pair the equal off-diagonal entries
and apply $2|z_\nu z_\mu|\le|z_\nu|^2+|z_\mu|^2$ to its quadratic
form. Its absolute Rayleigh quotient is bounded by the maximal row sum.

The image of a univariate polynomial of degree at most $N$ under CAI24
has total degree at most $N$. Multiplication by $\sum_i y_i$ has total
degree at most $N+1$, so no output is lost by the compression. Therefore
\[
 \|yP\|_{m_{h,k}}\le C_{\rm mult}(h)(2N+r_k)\|P\|_{m_{h,k}},\qquad
 \boxed{\frac{\omega_{N+1,k}^{\rm ar}}{\omega_{N,k}^{\rm ar}}
         \le C_{\rm mult}(h)^2(2N+r_k)^2}\quad(k\ge3,N\ge0).
 \tag{CAI25}
\]
For the last step $yQ_N$ is an actual degree-$N+1$ monic competitor;
its orthogonal monic minimum has no greater norm. The $S$-monic phases
have modulus one, so precisely the original norms occur. This proof
uses the full total-degree space, not independent degree-$N$ cutoffs
in each factor, and yields $N+k$ rather than $kN$.

\subsection{Both original contraction penalties on their complete window}

Return now to the stipulated packet in CAI1. By CAI11--12 and
$1-d_N\le1$,
\[
 d_n\le\frac{b_n^{\rm rec}}{b_n^{\rm rec}+L_{h,k}^2}
                         \quad(q\le n\le2q).
\]
For $n=q$ this is the endpoint formula; for $n\ge q+1$ it follows
from the formula at $N=n-1$. CAI25 bounds all these recurrence ratios by
\[
 B_{h,k}=C_{\rm mult}(h)^2(4q+r_k-2)^2,\qquad
 z_{h,k}=\frac{B_{h,k}}{L_{h,k}^2}
   =\frac{4C_{\rm mult}(h)^2(4+(r_k-2)/q)^2}{\delta^2(k+1)^2}.
 \tag{CAI26}
\]
The map $b\mapsto b/(b+L^2)$ is increasing for $b>0$, whence
\[
 0<d_n\le\frac{z_{h,k}}{1+z_{h,k}},\qquad
 \frac1{1+z_{h,k}}\le1-d_n<1,\qquad q\le n\le2q.\tag{CAI27}
\]
It is the complementary factors $1-d_n$, not $d_n$ themselves, that
approach one. In CAI13 the coefficient sum in $\mathcal P_-$ is
$1+2(q-1)=2q-1$, and that in $\mathcal P_+$ is $2+2(q-1)=2q$.
Thus the exact finite bound is
\[
 \boxed{0\le\mathcal P_-+\mathcal P_+
      \le\Pi_{h,k}:=(4q-1)\log(1+z_{h,k}).}\tag{CAI28}
\]
Since $q=e_k(k+1)^2$ and $r_k\le k/3$, $z_{h,k}=O_h(k^{-2})$ and
$\log(1+z)\le z$ show that $\Pi_{h,k}=O_h(q/k^2)=O_h(e_k)$.
This is $O_h(1)$ for $m=1$, and $O_h(k)=o(q)$ for any fixed $m>1$.
No early-$k$ smallness assumption was used in CAI27--28.
At an interior degree $N\ge q$, writing $\tau_N=-\log d_{N+1}$ gives
the exact pointwise comparison
$\epsilon_N^2/(b_{N+1}^{\rm rec}e^{\tau_N})=(1-d_N)(1-d_{N+1})$,
whose logarithm lies in $[-2\log(1+z_{h,k}),0]$ on the window.
At $N=q-1$ the one remaining factor is $1-d_q$.

\subsection{The four-norm factorial window and its finite remainder}

Minimize CAI18 over all monic polynomials, retaining the literal
degree-$n$ coefficient one. It gives
$(a_k/D_n)\gamma_n\le\omega_{n,k}^{\rm ar}\le A_k\gamma_n$.
Define
\[
\begin{gathered}
 \Lambda_{k,n}=\log(A_kD_n/a_k),\quad
 E_{\rm lo}=\Lambda_{k,q-1}+\Lambda_{k,q},\quad
 E_{\rm hi}=\Lambda_{k,2q-1}+\Lambda_{k,2q},\\
 \mathcal W_q^\sigma=\log\frac{\gamma_{2q-1}\gamma_{2q}}
                                     {\gamma_{q-1}\gamma_q},\qquad
 -E_{\rm hi}\le\mathcal W_k^{\rm ar}-\mathcal W_q^\sigma\le E_{\rm lo}.
\end{gathered}\tag{CAI29}
\]
For the lower bound put both numerator norms at their lower bounds
and both denominator norms at their upper bounds; for the upper bound
reverse those choices. The common $a_k,A_k$ terms are thereby counted
exactly twice. This comparison occurs \emph{after} the window has
telescoped, not once at each of its $q$ steps. More explicitly, the
original AMT22 constants give
\[
\begin{gathered}
 \mathfrak s_h=\alpha+\log(M_\alpha/\vartheta_h),\quad
 \mathfrak c_h=-\log M_\alpha+3\log\vartheta_h-3\alpha+
       \log\frac{C_h^{\rm tilt}C_\Gamma2^{B_h/2}}{c_hc_\Gamma},\\
 \log(A_k/a_k)=k\mathfrak s_h+(p+1)\log k+B_h\log(k-2)+\mathfrak c_h.
\end{gathered}\tag{CAI30}
\]
Since $h$ is fixed and $\log D_n=O_h(1+\log(n+1))$, the two errors
in CAI29 are $O_h(k+\log q)$.

The identity $(1/2)_n=(2n)!/(4^nn!)$ gives
$\gamma_n=\sqrt{2\pi}(2n)!/4^n$. Cancelling these four literal norms,
including their masses, gives
\[
 \mathcal W_q^\sigma
 =2\log\frac{(4q)!}{(2q)!}-4q\log2
                       +\log\frac{2q-1}{2(4q-1)}.\tag{CAI31}
\]
Here the endpoint factor results from
$(4q-2)!=(4q)!/[(4q)(4q-1)]$ and
$(2q-2)!=(2q)!/[(2q)(2q-1)]$; it is not discarded.

For a fully finite enclosure put
$C(q)=4q\log q+(8\log2-4)q$ and
$\eta_q=1/(16q)+1/(4q-2)$. For $f(x)=\log x$, integration by parts
twice on each unit interval $[j,j+1]$ gives
\[
 \frac{f(j)+f(j+1)}2-\int_j^{j+1}f(x)\,dx
 =\frac12\int_j^{j+1}(x-j)(j+1-x)f''(x)\,dx.
\]
Since $f''(x)=-x^{-2}$ and $(x-j)(j+1-x)/2\le1/8$, summing from
$j=2q$ to $4q-1$ bounds the composite trapezoidal error by
$-\frac18\int_{2q}^{4q}x^{-2}dx=-1/(32q)\le\xi_q\le0$.
The factorial sum is the trapezoidal sum plus
$[f(4q)-f(2q)]/2$. Substitution in CAI31 and exact integration give
\[
 \mathcal W_q^\sigma=C(q)-\log2+2\xi_q+
                                  \log(1-1/(4q-1)).
\]
Finally $-x/(1-x)\le\log(1-x)\le0$ for $0<x<1$, obtained by
integrating $-1/(1-x)$, proves for every integer $q\ge1$
\[
 C(q)-\log2-\eta_q\le\mathcal W_q^\sigma\le C(q)-\log2.
 \tag{CAI32}
\]

\subsection{Complete finite interval, action mean and exact receiving scope}

Combining CAI15, CAI28 and CAI29 proves the entirely specified interval
\[
 \boxed{\mathcal B_k^{\rm ar}+\mathcal W_q^\sigma-E_{\rm hi}-\Pi_{h,k}
       \le J_k^{\rm action}
       \le\mathcal B_k^{\rm ar}+\mathcal W_q^\sigma+E_{\rm lo}.}
 \tag{CAI33}
\]
In particular CAI32 gives
\[
\begin{aligned}
 \mathcal B_k^{\rm ar}+C(q)-\log2-\eta_q-E_{\rm hi}-\Pi_{h,k}
     &\le J_k^{\rm action},\\
 J_k^{\rm action}&\le\mathcal B_k^{\rm ar}+C(q)-\log2+E_{\rm lo}.
\end{aligned}\tag{CAI34}
\]
For fixed $h$ and its complete multiplicity, $e_k=O_h(k)$ and
$q=e_k(k+1)^2$. All noncentral terms in CAI34 are
$O_h(k+\log q)=o(q)$. Therefore
\[
 \boxed{J_k^{\rm action}=\mathcal B_k^{\rm ar}
     +4q\log q+(8\log2-4)q+O_h(k+\log q).}\tag{CAI35}
\]
The finer finite centre in CAI34 retains $-\log2$ and every error term.
The coefficient $8\log2-4$ is a coefficient of
$J_k^{\rm action}-\mathcal B_k^{\rm ar}$, not a coefficient of
$J_k^{\rm action}$ after replacing $\mathcal B_k^{\rm ar}$ by only its
leading asymptotic. The exact signed correction
$\mathcal B_k^{\rm ar}=\mathcal B_k^\sigma+\delta_k^\sigma$ remains
unevaluated at scale $q$. Nothing in CAI35 assigns either individual
kernel/residual allocation, or the absolute order-one long loss.

The original weighted arithmetic geometric mean is exactly
\[
 \mathcal G_k=\prod_{N=q-1}^{2q-1}\epsilon_N^{w_N/(2q)}
             =\exp(J_k^{\rm action}/(4q)).
\]
Its finite bounds are the exponentials of CAI34 divided by $4q$.
Consequently, since $(k+\log q)/q=O_h(1/k)$ for fixed $m$,
\[
 \boxed{\frac{\mathcal G_k}{q\exp(\mathcal B_k^{\rm ar}/(4q))}
          =\frac4{\exp(1)}(1+O_h(k^{-1}))\longrightarrow\frac4{\exp(1)}.}
 \tag{CAI36}
\]
This uses $e^x=1+O(x)$ as $x\to0$, which follows directly by integrating
the exponential derivative on a bounded interval. It also proves
$(J_k^{\rm action}-\mathcal B_k^{\rm ar})/(kq)\to0$ since
$\log q=O_h(\log k)$.

The exact exterior receiver is unchanged:
\[
 4q\log L_{h,k}\le J_k^{\rm action},\qquad
 L_{h,k}\le\min_{q-1\le N\le2q-1}\epsilon_N\le\mathcal G_k,
 \qquad
 \frac{\mathcal G_k}{L_{h,k}}=
 \frac8{\exp(1)\delta(k+1)}\exp(\mathcal B_k^{\rm ar}/(4q))
                              (1+O_h(k^{-1})).\tag{CAI37}
\]
The first two inequalities follow term by term from CAI11 and the
weights' sum $2q$; a minimum is bounded by, not identified with, the
weighted mean. In the original simple-quartet baseline domain alone,
AMT40 and the complete BSL/HBL source proof give
$\mathcal B_k^{\rm ar}/q^2\to C_B>769/17010>0$. Substitution into
CAI35 proves $J_k^{\rm action}/q^2\to C_B$, and CAI37 then tends to
infinity exponentially in $q$. This last consequence uses that already
proved baseline only on its original $m=1$ domain; the new interval,
penalty bound and relative expansion retain every complete multiplicity.

For the existing original simple-quartet outer receiver, retain its
literal finite identity
$\Psi_k=\mathcal F_{\partial,k}^{(k)}-\Phi_{k,k}
+\mathcal H_k+\delta_k^\sigma+\mathfrak e_{k,k}$, with the original
$-e_{k,k}^-\le\mathfrak e_{k,k}\le e_{k,k}^+$. Substituting CAI15 gives
\[
 \Phi_{k,k}\le\mathcal F_{\partial,k}^{(k)}+\mathcal H_k
       +\delta_k^\sigma+e_{k,k}^++C_k^{\rm budget}-4q\log L_{h,k}.
 \tag{CAI38}
\]
This is exactly OMR26's arithmetic backward receiver. Every signed low
metric and error remains where its original identity puts it. CAI33--34
can be substituted for the \emph{complete} action; they do not justify
independently improving one summand while holding its correlated budget
fixed. No failure of the broader programme or exclusion at other
admitted degrees follows from the size of this particular mean.
